\section{Sugerencia de ingeniería: el SFT primero afianza las trayectorias de comportamiento}
En los escenarios de Agentic RL, el valor del SFT multiturno no es solo «entrenar primero un modelo que sepa hablar», sino afianzar primero las trayectorias de comportamiento básicas del uso de herramientas. Solo cuando el modelo ya ha aprendido a invocar herramientas correctamente en un diálogo de varios turnos, leer la retroalimentación y continuar generando, el RL posterior dejará de desperdiciar gran parte del presupuesto de muestreo en los errores de formato más elementales.

\begin{knowledgebox}{Por qué hacer primero SFT y después RL}
\begin{itemize}
  \item El SFT aprende primero la plantilla de comportamiento y reduce la dificultad de exploración del RL
  \item Los datos multiturno le enseñan al modelo a manejar la estructura temporal de los mensajes de herramienta
  \item Una buena política inicial eleva de manera notable la proporción efectiva de los rollouts posteriores
\end{itemize}
\end{knowledgebox}

\subsection{Resumen de la sección}
Para los sistemas agenticos, el SFT es el punto de partida de la alineación de comportamiento. Primero fija los patrones básicos de interacción, y después deja que el RL optimice la calidad de la política con un grano más fino.

\newpage
